\section{Tour of \taoceorg Repositories}
%%%%%%%%%%%%%%%%%%%%%%%%%%%%%%%%%%%%%%%%%%%%%%%%%%%%%%%%%%%%%%%%%%%%%%%%%%%%%%%%
\subsection{Notice}
\begin{frame}
    \taoce is continuously improving, and the following slides can be considered accurate as-of today.\\
   
    \bigicon{\faRefresh}

    To enhance performance \& experiences of \TAOCE, major changes in \taoceorg repositories may happen in the future.
\end{frame}

%%%%%%%%%%%%%%%%%%%%%%%%%%%%%%%%%%%%%%%%%%%%%%%%%%%%%%%%%%%%%%%%%%%%%%%%%%%%%%%%
\subsection{\taoceorg Applications \& Packages}
\begin{frame}[fragile=singleslide]
    \only<article>{%
        \taoceorg repositories are organized as submodules. Therefore, applications and packages are split into separate repositories.
        An application may need one to many repositories to be built, and may depends on packages.
    }
    \only<presentation>{%
        \taocerepo is a \alert{monorepo}: a single repository gathering all applications and libraries.
    }
    \bigicon{\faFolderOpen}

    Currently, \taocerepo loads \alert{34} repositories as Git submodules.

\begin{onlyenv}<article>
You can list all submodules with the following command:
\inputminted[fontsize=\scriptsize]{text}{./snippets/git-submodule.txt}
\end{onlyenv}

\begin{onlyenv}<presentation>
    \begin{columns}[T]
        \scriptsize
    \begin{column}{.5\textwidth}
    \begin{itemize}
        \item \taocerepo[tao-portal]
        \item \taocerepo[tao-deliver-fe]
        \item \taocerepo[tao-deliver-be]
        \item \taocerepo[tao-dynamic-query-api]
        \item \taocerepo[proctoring]
    \end{itemize}
    \end{column}
    \begin{column}{.5\textwidth}
    \begin{itemize}
        \item \taocerepo[tao-manual-scoring-be]
        \item \taocerepo[tao-manual-scoring-fe]
        \item \taocerepo[live-design-system]
        \item \taocerepo[read-aloud-client-fe]
        \item \dots
    \end{itemize}
    \end{column}
    \end{columns}
\end{onlyenv}
\end{frame}

%%%%%%%%%%%%%%%%%%%%%%%%%%%%%%%%%%%%%%%%%%%%%%%%%%%%%%%%%%%%%%%%%%%%%%%%%%%%%%%%
\subsection{\taocerepo Structure}
\begin{frame}[t,fragile]
\only<article>{%
    \taocerepo repository is structured as a \textit{monorepo} to gather all \taoce code, in a single file-tree.
}
\begin{dirtreex}[box={false},fontsize=\small]
    \dir{tao-ce}{}{
        \dir{apps}{}{
        \only<2->{
            \dir{portal}{}{ \dir{src}{\sffamily\taocerepo[tao-portal] submodule}{} }
            \dir{deliver}{}{
                \dir{src}{}{
                    \dir{fe}{\sffamily\taocerepo[tao-deliver-fe] submodule}{}
                    \dir{be}{\sffamily\taocerepo[tao-deliver-be] submodule}{}
            }}
            \dir{\dots}{\sffamily other \TAOCE apps}{}
        }}
        \dir{packages}{}{
        \only<3->{
            \dir{\dots}{\sffamily \TAOCE libraries}{}
        }}
        \only<4->{
        \dir{build}{}{}
        \dir{etc}{}{}
        \dir{docs}{}{}
        \dir{libexec}{}{}
        \dir{\dots}{}{}
        }
    }
\end{dirtreex}
\end{frame}

%%%%%%%%%%%%%%%%%%%%%%%%%%%%%%%%%%%%%%%%%%%%%%%%%%%%%%%%%%%%%%%%%%%%%%%%%%%%%%%%
\subsection{Build Resources}
\begin{frame}[t]
\only<article>{%
    The \texttt{build} directory contains resources separated in stages, related to artifacts we intend to build.
}
\begin{dirtreex}[box={false},fontsize=\small]
    \dir{tao-ce}{}{
        \dir{build}{}{
            \dir{crystal}{\sffamily Build \TAOCE from sources}{
                \dir{Containerfile.d}{\sffamily Containerfile fragments}{}
            }
            \dir{swift}{\sffamily Buid \TAOCE images}{}
            \dir{docs}{\sffamily Build \href{https://tao-ce.github.io/tao-ce}{TAO CE docs}}{}
            \dir{\dots}{}{}
        }
        \file{Containerfile}{\sffamily Base container image}
        \file{Taskfile.yml}{\sffamily Build and operation tasks}
        \file{\dots}{}
    }
\end{dirtreex}
\end{frame}

%%%%%%%%%%%%%%%%%%%%%%%%%%%%%%%%%%%%%%%%%%%%%%%%%%%%%%%%%%%%%%%%%%%%%%%%%%%%%%%%
\subsection{Runtime Resources}
\begin{frame}[t]
\begin{dirtreex}[box={false},fontsize=\small]
    \dir{tao-ce}{}{
        \dir{etc}{}{
        \only<2->{
            \dir{tao-ce}{}{
                \dir{config}{}{
                    \file{Caddyfile}{\sffamily Main endpoint configuration}
                    \file{tao.yaml}{\sffamily Main \TAOCE configuration}
                }
            }
        }}
        \dir{libexec}{}{
            \only<3->{
            \dir{init/portal}{\sffamily Default Portal users}{}
            \dir{setup}{\sffamily Configuration templates}{}
            \dir{systemd}{\sffamily Systemd services}{}
            \dir{\dots}{}{}
        }}
    }
\end{dirtreex}
\end{frame}

%%%%%%%%%%%%%%%%%%%%%%%%%%%%%%%%%%%%%%%%%%%%%%%%%%%%%%%%%%%%%%%%%%%%%%%%%%%%%%%%
\subsection<article>{\taocerepo Documentation}
\begin{frame}<article>
    \bigicon{\faBook}
\begin{itemize}
    \item \url{https://tao-ce.github.io/tao-ce/}
    \item rebuilt automatically from \href{https://github.com/tao-ce/tao-ce/tree/main/docs}{\texttt{docs}}
\end{itemize}
\end{frame}

